% GENERATED FILE. Edit canonical lessons, then run npm run generate:books.
% canonical-source-hash: fnv1a64:56297a8daddb905c
% canonical-lessons: PT-C400-publicar, PT-C400-comentar, PT-C400-seguidor, PT-C400-reencaminhar, PT-C400-bloquear, PT-R400-revisao-399-400, PT-R400-revisao-final-400

\chapter{Nas redes sociais}
\label{ch:pt-nas-redes-sociais}
\hlchaptermodality{400}

\begin{chapteropening}
\textbf{By the end of this chapter:} \emph{I can use social media: post and comment, see my followers, forward messages and block someone.}

Complete the last lesson of chapter 400: I can use social media: post and comment, see my followers, forward messages and block someone.
\end{chapteropening}

\section[publicar]{publicar — to post, to publish}

\label{lesson:PT-C400-publicar}



\begin{quote}
Before the new one: say the Portuguese for the search engine, then the Portuguese for the web link.
\end{quote}

\subsection*{You'll want to know: publicar}
\textbf{publicar} — \textquotedblleft{}to post, to publish\textquotedblright{}.

\begin{quote}
Ela publicou as fotos da viagem.
\end{quote}

\begin{quote}
\emph{She posted the photos of the trip.}
\end{quote}

\emph{Publicar um livro} — to publish a book.

\begin{cousinweb}
\emph{publicar} is a learned borrowing from Latin \emph{pūblicāre} \textquotedblleft{}to make public\textquotedblright{} (EtymDB), on \emph{pūblicus} \textquotedblleft{}of the people, public\textquotedblright{} (Lewis \& Short).
\end{cousinweb}

\begin{grammarlens}[title={What you've built}]
One more word for telling, explaining and describing at B1.
\end{grammarlens}

\subsection*{Your turn}
\begin{itemize}
\raggedright
  \item \emph{Say it:} \emph{publicar}
  \item \emph{Say it:} \emph{publicar}, once more
  \item \emph{Say it:} a sentence of your own with \emph{publicar}
  \item \emph{Recall:} say the Portuguese for the search engine, then the Portuguese for the web link, then say \emph{publicar} again
\end{itemize}

\begin{tcolorbox}[breakable,colback=teal!4,colframe=teal!35!black,title={Before you move on}]
What does publicar mean? (\textquotedblleft{}To post, to publish\textquotedblright{}.) Say it once more.
\end{tcolorbox}
\section[comentar]{comentar — to comment (on)}

\label{lesson:PT-C400-comentar}



\begin{quote}
Before the new one: say the Portuguese for the web link, then the Portuguese for to post.
\end{quote}

\subsection*{You'll want to know: comentar}
\textbf{comentar} — \textquotedblleft{}to comment (on)\textquotedblright{}.

\begin{quote}
Toda a gente comentou a fotografia.
\end{quote}

\begin{quote}
\emph{Everyone commented on the photo.}
\end{quote}

\emph{Comentar alguma coisa} — no preposition needed.

\begin{cousinweb}
\emph{comentar} is a learned borrowing from Latin \emph{commentārī} \textquotedblleft{}to think over, to make notes on\textquotedblright{} (EtymDB), a verb formed on \emph{comminīscī} \textquotedblleft{}to devise, to recall\textquotedblright{} (Lewis \& Short).
\end{cousinweb}

\begin{grammarlens}[title={What you've built}]
One more word for telling, explaining and describing at B1.
\end{grammarlens}

\subsection*{Your turn}
\begin{itemize}
\raggedright
  \item \emph{Say it:} \emph{comentar}
  \item \emph{Say it:} \emph{comentar}, once more
  \item \emph{Say it:} a sentence of your own with \emph{comentar}
  \item \emph{Recall:} say the Portuguese for the web link, then the Portuguese for to post, then say \emph{comentar} again
\end{itemize}

\begin{tcolorbox}[breakable,colback=teal!4,colframe=teal!35!black,title={Before you move on}]
What does comentar mean? (\textquotedblleft{}To comment (on)\textquotedblright{}.) Say it once more.
\end{tcolorbox}
\section[o seguidor]{o seguidor — the follower}

\label{lesson:PT-C400-seguidor}



\begin{quote}
Before the new one: say the Portuguese for to post, then the Portuguese for to comment.
\end{quote}

\subsection*{You'll want to know: o seguidor}
\textbf{o seguidor} — \textquotedblleft{}the follower\textquotedblright{}.

\begin{quote}
Ela tem dez mil seguidores.
\end{quote}

\begin{quote}
\emph{She has ten thousand followers.}
\end{quote}

Feminine \emph{a seguidora}.

\begin{cousinweb}
\emph{seguidor} is built on \emph{seguir} \textquotedblleft{}to follow\textquotedblright{}, inherited from Latin \emph{sequī} \textquotedblleft{}to follow\textquotedblright{} (Lewis \& Short), with \emph{-dor}.
\end{cousinweb}

\begin{grammarlens}[title={What you've built}]
One more word for telling, explaining and describing at B1.
\end{grammarlens}

\subsection*{Your turn}
\begin{itemize}
\raggedright
  \item \emph{Say it:} \emph{o seguidor}
  \item \emph{Say it:} \emph{o seguidor}, once more
  \item \emph{Say it:} a sentence of your own with \emph{o seguidor}
  \item \emph{Recall:} say the Portuguese for to post, then the Portuguese for to comment, then say \emph{o seguidor} again
\end{itemize}

\begin{tcolorbox}[breakable,colback=teal!4,colframe=teal!35!black,title={Before you move on}]
What does o seguidor mean? (\textquotedblleft{}The follower\textquotedblright{}.) Say it once more.
\end{tcolorbox}
\section[reencaminhar]{reencaminhar — to forward (a message)}

\label{lesson:PT-C400-reencaminhar}



\begin{quote}
Before the new one: say the Portuguese for to comment, then the Portuguese for the follower.
\end{quote}

\subsection*{You'll want to know: reencaminhar}
\textbf{reencaminhar} — \textquotedblleft{}to forward (a message)\textquotedblright{}.

\begin{quote}
Reencaminhei a mensagem para a minha irmã.
\end{quote}

\begin{quote}
\emph{I forwarded the message to my sister.}
\end{quote}

European Portuguese; in Brazil \emph{encaminhar}.

\begin{cousinweb}
\emph{reencaminhar} is built from \emph{re-}, \emph{en-} and \emph{caminho} \textquotedblleft{}way\textquotedblright{}, which is inherited from Late Latin \emph{cammīnus} \textquotedblleft{}way\textquotedblright{} (EtymDB), usually traced to a Celtic word.
\end{cousinweb}

\begin{grammarlens}[title={What you've built}]
One more word for telling, explaining and describing at B1.
\end{grammarlens}

\subsection*{Your turn}
\begin{itemize}
\raggedright
  \item \emph{Say it:} \emph{reencaminhar}
  \item \emph{Say it:} \emph{reencaminhar}, once more
  \item \emph{Say it:} a sentence of your own with \emph{reencaminhar}
  \item \emph{Recall:} say the Portuguese for to comment, then the Portuguese for the follower, then say \emph{reencaminhar} again
\end{itemize}

\begin{tcolorbox}[breakable,colback=teal!4,colframe=teal!35!black,title={Before you move on}]
What does reencaminhar mean? (\textquotedblleft{}To forward (a message)\textquotedblright{}.) Say it once more.
\end{tcolorbox}
\section[bloquear]{bloquear — to block; to freeze (a screen)}

\label{lesson:PT-C400-bloquear}



\begin{quote}
Before the new one: say the Portuguese for the follower, then the Portuguese for to forward.
\end{quote}

\subsection*{You'll want to know: bloquear}
\textbf{bloquear} — \textquotedblleft{}to block; to freeze (a screen)\textquotedblright{}.

\begin{quote}
Bloqueei um utilizador que me insultava.
\end{quote}

\begin{quote}
\emph{I blocked a user who was insulting me.}
\end{quote}

\emph{O ecrã bloqueou} — the screen froze.

\begin{cousinweb}
\emph{bloquear} is usually traced to French \emph{bloquer} \textquotedblleft{}to block\textquotedblright{}, the source EtymDB gives for its Spanish cousin \emph{bloquear}; \emph{bloquer} is built on \emph{bloc}, which came through Old French \emph{bloc} \textquotedblleft{}log, block\textquotedblright{} from Middle Dutch \emph{blok} \textquotedblleft{}a tree trunk\textquotedblright{} (EtymDB).
\end{cousinweb}

\begin{grammarlens}[title={What you've built}]
One more word for telling, explaining and describing at B1.
\end{grammarlens}

\subsection*{Your turn}
\begin{itemize}
\raggedright
  \item \emph{Say it:} \emph{bloquear}
  \item \emph{Say it:} \emph{bloquear}, once more
  \item \emph{Say it:} a sentence of your own with \emph{bloquear}
  \item \emph{Recall:} say the Portuguese for the follower, then the Portuguese for to forward, then say \emph{bloquear} again
\end{itemize}

\begin{tcolorbox}[breakable,colback=teal!4,colframe=teal!35!black,title={Before you move on}]
What does bloquear mean? (\textquotedblleft{}To block; to freeze (a screen)\textquotedblright{}.) Say it once more.
\end{tcolorbox}
\section[(dialogue)]{Condições e regras}

\label{lesson:PT-R400-revisao-399-400}



\begin{quote}
Say the Portuguese for the network and to browse.
\end{quote}

\subsection*{Your turn}
Say these aloud:

\begin{itemize}
\raggedright
  \item the network, to browse, to search, the search engine, the web link
  \item to post, to comment, the follower, to forward, to block
\end{itemize}

\begin{tcolorbox}[breakable,colback=teal!4,colframe=teal!35!black,title={Before you move on}]
The network? (\textbf{a rede}.) To browse? (\textbf{navegar}.) To search? (\textbf{pesquisar}.) The search engine? (\textbf{o motor de busca}.) The web link? (\textbf{a hiperligação}.) To post? (\textbf{publicar}.) To comment? (\textbf{comentar}.) The follower? (\textbf{o seguidor}.) To forward? (\textbf{reencaminhar}.) To block? (\textbf{bloquear}.)
\end{tcolorbox}
\section[(dialogue)]{Nas redes sociais — once more}

\label{lesson:PT-R400-revisao-final-400}



\begin{quote}
Say the Portuguese for to post and to comment.
\end{quote}

\subsection*{Your turn}
\begin{itemize}
\raggedright
  \item \emph{Say it:} to post, to comment, the follower, to forward, to block
\end{itemize}

\begin{tcolorbox}[breakable,colback=teal!4,colframe=teal!35!black,title={Before you move on}]
To post? (\textbf{publicar}.) To comment? (\textbf{comentar}.) The follower? (\textbf{o seguidor}.) To forward? (\textbf{reencaminhar}.) To block? (\textbf{bloquear}.)
\end{tcolorbox}
